% %%%%%%%%%%%%%%%%%%%%%%%%%%%%%%%%%%%%%%%%%%%%%%%%%%%%%%%%%%%%%%%%%%%%%
% %% State-of-the-Art %%
% SortCache
% %%%%%%%%%%%%%%%%%%%%%%%%%%%%%%%%%%%%%%%%%%%%%%%%%%%%%%%%%%%%%%%%%%%%%

\subsection{SortCache}
\label{sec:soa:sortcache}

The authors of SortCache~\cite{Srikanth2021_SortCache} note that modern processors are able to load an entire cache line in parallel but in many sparse applications less than 50\% of the data within a line is used, even with hand-tuned code. They propose a light SpGEMM accelerator which offloads the \emph{reduction} step, based on the \emph{Vectorized Binary Search Tree} (VBST), a data-structure from their earlier work~\cite{Srikanth2017_SuperStrider,Srikanth2019_MetaStrider}.
The VBST is a standard binary search tree where the size of node is a multiple of the cache block size and contains a sorted vector of $K$ $(key, value)$ pairs. Associated with each node is a \emph{pivot}; the pivots, with left and right child pointers, form a binary tree, as illustrated in Figure~\ref{fig:soa:sortcache_vbst}.

% Location of Figure in Accelerators.tex

Using the VBST, SortCache focuses on accelerating the \emph{scatter} updates to the output matrix, which are performed directly in the cache. The authors augment the processor with a single new instruction, \mbox{\emph{RED Key, Value}}. After $K$ updates have been issued, the hardware launches a \emph{reduction} in the SortCache, where the $K$ $(key, value)$ pairs are inserted into the VBST. This is done by traversing the tree from the root to the leaves. When duplicate keys are found, the \emph{reduction} operation is performed in the cache. When new keys are found, they are inserted, which can result in a temporary, new node with up to $2K$ $(key, value)$ pairs. This temporary node is partitioned, with one half having exactly $K$ entries. The insertion procedure may require visiting each level of the tree at least once ($\mathcal{O}(log(n)$) and it does not ensure the uniqueness of the keys, so at the end, an additional full traversal of the tree is required to remove duplicates.

The nodes closest to the root of the tree are accessed most frequently and are thus mapped to the caches closest to the processor. The authors propose to re-purpose cache tracking hardware (e.g. TAG values, LRU counters) to store the additional metadata (pivots, pointers). One of the difficulties with SortCache is that, in some cases, a large fraction of the VBST nodes are underutilized.